%%%%%
%
\section{Transformation du temps GPS en temps UTC}

\itemstyle{Op\'erateur~:} \operfmt{GPS\_to\_UTC}
\index{GPS\_to\_UTC}
\label{OP-GPS_to_UTC}

%%%%%OBJET%%%%
\itemstyle{Objet~:}
Le temps GPS est compt� en secondes depuis le d�but de
la semaine, dans le syst�me GPS (d�cal� d'une valeur
constante avec le TAI et d'une valeur variable avec l'UTC).
Cet op�rateur
convertit des temps exprim�s dans le syst�me GPS en temps UTC
(compt� en secondes depuis 0h00). 
Le d�calage GPS/UTC doit �tre explicitement donn� en argument.

%%%%%ARGUMENTS%%%%
\bigskip
%%%%%ENTREES%%%%
\begin{tabarg}
\multicolumn{3}{l}{\titretab {Entr\'ees}}\\
\hline
$Temps_{GPS}$ & Vecteur, \texttt{float} & Temps dans le syst�me GPS (second)\\
$Decal$ & Coeff., {\tt Int\_32         } & D�calage GPS/UTC � la date consid�r�e (positif) (second)\\
\end{tabarg}

\medskip
%%%%%SORTIES%%%%
\begin{tabarg}
\multicolumn{3}{l}{\titretab {Sorties}}\\
\hline
$Temps_{UTC}$ & Vecteur, \texttt{double} & Temps, exprim� dans le syst�me UTC (second)\\
\end{tabarg}


%%%%%FORMULE%%%%
\itemstyle{Formule}

%%GPS_to_UTC
  Calcul d'un offset sur le premier $Temps_{GPS}$ non-nul :
\[
offset = 86400 . [(Temps_{GPS}-Decal) / 86400]
\]
Puis :
\[
   Temps_{UTC} = Temps_{GPS} - Decal - offset
\]


%%%%%COMMENTAIRE%%%%
\itemstyle{Commentaire~:}
Si un vecteur des diff�rences GPS-UTC est disponible, il est plus simple 
d'utiliser \hyperref[OP-GPS_to_UTC_auto]{\operfmt{GPS\_to\_UTC\_auto}}.

%%%%%AUTEUR%%%%
\itemstyle{Auteur~:}
CNRM/GMEI/TRAMM

%%%%%REFERENCE%%%%
\itemstyle{R\'ef\'erences~:}
\textit{GPS, theory and practice}, par B. Hofmann-Wellenhof, H. Lichtenegger, and J. Collins \cite{Hofmann-GPS}
